\chapter{From Boundary Conditions to the Open-String Spectrum}

\lectureoverview{Our aim is to quantize an open string and decide whether its
states deserve to be called particles. The route is elementary but exacting.
We first prepare the calculus and Fourier analysis of functions on
\(0\leq\sigma\leq\pi\). We then replace pointlikeness by a spectral criterion
for particle behavior, review the infinite-momentum description in which the
transverse dynamics become nonrelativistic, and build the string from a chain
of light masses and stiff springs. The continuum string is a wave system with
free endpoints. Its Neumann modes decouple into harmonic oscillators with
frequencies \(1,2,3,\ldots\); quantization therefore produces a discrete,
gapped spectrum. String interactions finally force us beyond open strings:
an open string can join its own endpoints and become closed.}

\section{Calculus on the open-string interval}

The calculation will repeatedly turn discrete sums into integrals and
nearest-neighbor differences into derivatives. It is better to put those
facts in place before they interrupt the physics.

Let \(x(\sigma)\) be a function on

\begin{equation}
  0\leq \sigma\leq \pi .
\end{equation}

The notation anticipates string theory: \(\sigma\) labels position along the
string, while \(x\) will become one of its transverse coordinates. Divide the
interval into \(N\) pieces,

\begin{equation}
  \Delta\sigma=\frac{\pi}{N},
  \qquad
  \sigma_i=i\,\Delta\sigma,
  \qquad
  x_i=x(\sigma_i).
  \label{eq:l02-lattice}
\end{equation}

For a smooth limiting function,

\begin{equation}
  \Delta x_i\equiv x_i-x_{i-1}
  =
  \frac{\partial x}{\partial\sigma}\,\Delta\sigma
  +O\!\left((\Delta\sigma)^2\right),
  \label{eq:l02-difference}
\end{equation}

and the Riemann sum becomes

\begin{equation}
  \sum_i x_i\,\Delta\sigma
  \longrightarrow
  \int_0^\pi x(\sigma)\,d\sigma .
  \label{eq:l02-riemann}
\end{equation}

\begin{figure}[H]
  \centering
  \includegraphics[width=0.78\textwidth]{lecture_02_figure_01.png}
  \caption{The finite-difference and Riemann-sum rules kept on the board for
  the later continuum limit.}
  \lectureframe{00:04:20}
  \label{fig:l02-calculus-board}
\end{figure}

These are definitions familiar from calculus. Their importance here is in
the scaling: \(\Delta x_i\) brings one power of \(\Delta\sigma\), whereas a
sum brings one inverse power when it becomes an integral.

\subsection{Dirichlet and Neumann endpoints}

Fourier analysis on a finite interval begins with a choice at the endpoints.
Two choices will recur.

\begin{align}
  \text{Dirichlet:}\qquad
  &x(0)=x(\pi)=0,
  \label{eq:l02-dirichlet}\\
  \text{Neumann:}\qquad
  &\left.\partial_\sigma x\right|_{\sigma=0}
   =
   \left.\partial_\sigma x\right|_{\sigma=\pi}
   =0.
  \label{eq:l02-neumann}
\end{align}

A violin string held at both ends gives the Dirichlet picture: its
displacement vanishes at the boundaries. A Neumann function need not vanish
there; it meets each boundary with zero slope. The two characteristic shapes
are visible together on the board.

\begin{figure}[H]
  \centering
  \begin{minipage}[t]{0.48\textwidth}
    \centering
    \includegraphics[width=\textwidth]{lecture_02_figure_02.png}
    \par\smallskip\footnotesize (a) Pinned and flat endpoint profiles.
  \end{minipage}\hfill
  \begin{minipage}[t]{0.48\textwidth}
    \centering
    \includegraphics[width=\textwidth]{lecture_02_figure_03.png}
    \par\smallskip\footnotesize (b) Light rings sliding on vertical poles.
  \end{minipage}
  \caption{Boundary conditions as shapes and as mechanics. The sliding-ring
  construction gives a free endpoint an unambiguous transverse motion.}
  \lectureframe{00:08:50 and 00:12:15}
  \label{fig:l02-boundaries-board}
\end{figure}

\begin{figure}[H]
  \centering
  \begin{tikzpicture}[>=Latex,thick,scale=0.92]
    \begin{scope}
      \draw[->] (0,0) -- (4.4,0) node[right] {\(\sigma\)};
      \draw[->] (0,-1.25) -- (0,1.35) node[above] {\(x\)};
      \draw[blue!70!black]
        (0,0) .. controls (0.8,1.1) and (1.4,1.0) .. (2.0,0)
        .. controls (2.7,-1.0) and (3.4,-0.8) .. (4.1,0);
      \fill (0,0) circle (0.06);
      \fill (4.1,0) circle (0.06);
      \node at (2.05,-1.65) {Dirichlet: \(x=0\)};
    \end{scope}
    \begin{scope}[xshift=5.8cm]
      \draw[->] (0,0) -- (4.4,0) node[right] {\(\sigma\)};
      \draw[->] (0,-1.25) -- (0,1.35) node[above] {\(x\)};
      \draw[orange!85!black]
        (0,0.65) .. controls (0.8,0.65) and (1.1,-0.9) .. (2.0,-0.9)
        .. controls (2.9,-0.9) and (3.2,0.55) .. (4.1,0.55);
      \draw[orange!85!black] (0,0.65) -- (0.45,0.65);
      \draw[orange!85!black] (3.65,0.55) -- (4.1,0.55);
      \node at (2.05,-1.65) {Neumann: \(\partial_\sigma x=0\)};
    \end{scope}
  \end{tikzpicture}
  \caption{The value is fixed for Dirichlet data; the slope is fixed for
  Neumann data.}
  \label{fig:l02-boundaries-redraw}
\end{figure}

\begin{classroomqa}
\audiencequestion{Does an open or closed organ pipe supply a clean example of
one of these boundary conditions?\lecturetimestamp{00:09:19}}
\lecturerresponse{It can, but one must first say which acoustic variable is
being expanded. Pressure perturbation and longitudinal displacement obey
complementary idealized conditions at an open end. Rather than settle that
side question at the board, use the sliding-ring model: nearly massless rings
can move vertically on the poles, so the string endpoints are free and the
Neumann picture is direct.}
\end{classroomqa}

\begin{classroomqa}
\audiencequestion{Can the graph rise vertically and contain a sharp
jump?\lecturetimestamp{00:09:53}}
\lecturerresponse{Fourier analysis can represent piecewise-continuous
functions with isolated jumps. A physical string is more restrictive. A
jump would disconnect neighboring points; in the continuum energy derived
below it also makes \(\partial_\sigma x\) singular. For the string shapes of
interest we therefore keep a continuous, differentiable profile.}
\end{classroomqa}

\begin{classroomqa}
\audiencequestion{Why should the open end of a pipe, or any free end, satisfy
a boundary condition at all?\lecturetimestamp{00:11:14}}
\lecturerresponse{The acoustic example depends on the chosen field and was
left unresolved in the discussion. The mechanical example does not. A very
light ring is free to slide but is still attached to the final string
segment. We shall derive its condition from Newton's law rather than assume
it: finite endpoint acceleration forces the slope to vanish.}
\end{classroomqa}

\subsection{The Fourier bases}

Once the endpoint class is fixed, the corresponding Fourier basis is fixed.
On \(0\leq\sigma\leq\pi\),

\begin{align}
  x_D(\sigma)
  &=
  \sum_{n=1}^{\infty}x_n\sin(n\sigma),
  \label{eq:l02-dirichlet-series}\\
  x_N(\sigma)
  &=
  \sum_{n=0}^{\infty}x_n\cos(n\sigma).
  \label{eq:l02-neumann-series}
\end{align}

Every sine in \eqref{eq:l02-dirichlet-series} vanishes at both endpoints.
Every cosine in \eqref{eq:l02-neumann-series} has zero endpoint derivative.
The lower limits differ for a reason:
\(\sin(0\sigma)=0\) contributes nothing, while
\(\cos(0\sigma)=1\) is the constant, or zero, mode. It records the average
value of a Neumann function.

\begin{figure}[H]
  \centering
  \begin{minipage}[t]{0.48\textwidth}
    \centering
    \includegraphics[width=\textwidth]{lecture_02_figure_04.png}
    \par\smallskip\footnotesize (a) Sine and cosine expansions.
  \end{minipage}\hfill
  \begin{minipage}[t]{0.48\textwidth}
    \centering
    \includegraphics[width=\textwidth]{lecture_02_figure_05.png}
    \par\smallskip\footnotesize (b) Orthogonality and the exceptional
    constant mode.
  \end{minipage}
  \caption{The Fourier facts used later to diagonalize the string.}
  \lectureframe{00:16:15 and 00:20:30}
  \label{fig:l02-fourier-board}
\end{figure}

Differentiation exchanges the two classes:

\begin{equation}
  \partial_\sigma\sin(n\sigma)=n\cos(n\sigma),
  \qquad
  \partial_\sigma\cos(n\sigma)=-n\sin(n\sigma).
  \label{eq:l02-basis-derivatives}
\end{equation}

Thus the derivative of a Dirichlet series has Neumann behavior and the
derivative of a Neumann series has Dirichlet behavior. Orthogonality then
removes cross-terms between distinct modes:

\begin{align}
  \int_0^\pi\cos(n\sigma)\cos(m\sigma)\,d\sigma
  &=
  \frac{\pi}{2}\delta_{nm},
  && n,m\geq1,
  \label{eq:l02-cos-orthogonality}\\
  \int_0^\pi\sin(n\sigma)\sin(m\sigma)\,d\sigma
  &=
  \frac{\pi}{2}\delta_{nm},
  && n,m\geq1,
  \label{eq:l02-sin-orthogonality}\\
  \int_0^\pi 1\,d\sigma&=\pi .
  \label{eq:l02-zero-norm}
\end{align}

The last line is the exceptional \(n=m=0\) cosine norm. A constant is not
normalized like a nonzero harmonic, and that factor of two will distinguish
the string's center-of-mass mode from its internal modes.

\subsection{One oscillator formula to remember}

After rescaling a coordinate so that its mechanical mass is one, a
Hooke-law oscillator has

\begin{align}
  E_{\mathrm{osc}}
  &=
  \frac12\dot q^{\,2}
  +\frac12\omega^2q^2,
  \label{eq:l02-oscillator-energy}\\
  L_{\mathrm{osc}}
  &=
  \frac12\dot q^{\,2}
  -\frac12\omega^2q^2.
  \label{eq:l02-oscillator-lagrangian}
\end{align}

\begin{figure}[H]
  \centering
  \includegraphics[width=0.68\textwidth]{lecture_02_figure_06.png}
  \caption{The oscillator energy, written so that the coefficient of
  \(q^2\) immediately identifies \(\omega^2\).}
  \lectureframe{00:24:00}
  \label{fig:l02-oscillator-board}
\end{figure}

Classically, \(\omega=\sqrt{\kappa/m}\); the coordinate rescaling has set
\(m=1\). Quantum mechanically,

\begin{equation}
  E_Q=\hbar\omega\left(Q+\frac12\right),
  \qquad Q=0,1,2,\ldots .
  \label{eq:l02-oscillator-spectrum}
\end{equation}

For level spacings we may subtract the common ground-state term and say that
each quantum costs \(\hbar\omega\). The string calculation will initially do
exactly that and postpone the total zero-point energy.

\begin{classroomqa}
\audiencequestion{How is the oscillator frequency read from the energy
formula?\lecturetimestamp{00:23:21}}
\lecturerresponse{Before rescaling it is
\(\omega=\sqrt{\kappa/m}\). With unit mass the potential term is
\(\frac12\omega^2q^2\), so its quadratic coefficient is the square of the
frequency.}
\end{classroomqa}

\section{What makes a composite system a particle?}

A string is extended and will be modeled by many coupled pieces. In what
sense can one of its quantum states nevertheless be a particle? Pointlikeness
cannot be the decisive criterion. Protons and neutrons are composite.
Electrons and photons are accompanied by quantum-field dressing even though
no intrinsic substructure has been resolved. A classical geometric point
surrounded by nothing is not the right image of a quantum excitation.

Nor do mass and center-of-mass position suffice. A sealed box of gas has
both. A cup of coffee has both. We do not normally call either one a
particle. Indivisibility also fails because a proton is divisible, and spin
does not distinguish microscopic particles from macroscopic bodies.

The useful criterion is spectral isolation. Put the rest energy on a vertical
axis and ask how much energy is required to reach another internal state. An
electron has no experimentally accessible tower of electron excitations just
above it. A proton does have resonances, but the first ones cost a few hundred
MeV, a substantial fraction of the proton mass. Those states are separated
enough to be individually identifiable.

A cup of coffee is the opposite. One can add an arbitrarily small amount of
energy to a long-wavelength sound mode, nudge a molecule, or produce a phonon.
There are enormously many states packed into a tiny interval. On laboratory
scales they form a continuum. This is the practical distinction between a
particle and \emph{mush}.

\begin{figure}[H]
  \centering
  \includegraphics[width=0.58\textwidth]{lecture_02_figure_07.png}
  \caption{The board's vertical energy axis during the comparison of an
  isolated particle level with the dense excitations of ordinary matter.}
  \lectureframe{00:31:10}
  \label{fig:l02-spectrum-board}
\end{figure}

\begin{figure}[H]
  \centering
  \begin{tikzpicture}[>=Latex,thick]
    \begin{scope}
      \draw[->] (0,0) -- (0,4.8) node[above] {energy};
      \draw[blue!70!black,line width=1.2pt] (-1.0,0.65) -- (1.0,0.65);
      \draw[blue!70!black,line width=1.2pt] (-1.0,3.65) -- (1.0,3.65);
      \node[right] at (1.05,0.65) {ground state};
      \node[right] at (1.05,3.65) {first internal excitation};
      \draw[<->] (0.65,0.75) -- (0.65,3.55)
        node[midway,right] {large gap};
      \node at (0,-0.55) {particle-like};
    \end{scope}
    \begin{scope}[xshift=7.0cm]
      \draw[->] (0,0) -- (0,4.8) node[above] {energy};
      \foreach \y in {0.55,0.69,0.82,0.94,1.05,1.15,1.24,1.32,1.39,1.45,
        1.51,1.57,1.63,1.70,1.78,1.87,1.97,2.09,2.22,2.36,2.51,2.67,2.84,
        3.02,3.21,3.41,3.62,3.84,4.07}
        \draw[orange!85!black] (-0.95,\y) -- (0.95,\y);
      \node[right] at (1.05,2.25) {dense phonon and internal states};
      \node at (0,-0.55) {macroscopic matter};
    \end{scope}
  \end{tikzpicture}
  \caption{Spectral isolation, not literal pointlikeness, is the operational
  particle criterion used in the lecture.}
  \label{fig:l02-particle-redraw}
\end{figure}

This criterion turns the philosophical question into a calculation. Quantize
the string, find its internal energy levels, and inspect their spacing. If
the first excited string state lies near the Planck scale while the observed
ground state is light, the gap exceeds present accelerator energies by more
than twenty orders of magnitude. The string then behaves as an elementary
particle throughout the accessible regime.

\begin{classroomqa}
\audiencequestion{Is particle behavior merely an issue of observational
scale? If we zoom in far enough, do the levels cease to be
separate?\lecturetimestamp{00:33:33}}
\lecturerresponse{Zooming in may reveal that an object has many constituents,
but it does not erase the spectrum. The relevant comparison is between the
energy resolution of the experiment and the gap to the first internal
excitation. A composite object can remain particle-like when that gap is
enormous.}
\end{classroomqa}

\begin{classroomqa}
\audiencequestion{Why is an ordinary violin string not a
particle?\lecturetimestamp{00:34:37}}
\lecturerresponse{One vibrational quantum carries an immeasurably small
energy compared with the rest energy of the macroscopic string. Many quantum
states therefore fit into a tiny relative interval. It belongs on the dense
side of the spectral comparison, not the isolated side.}
\end{classroomqa}

\begin{classroomqa}
\audiencequestion{Are photons and electrons point particles if quantum fields
surround them with virtual fluctuations?\lecturetimestamp{00:35:40}}
\lecturerresponse{The classical word \emph{point} is doing too much work. The
physical electron and photon are dressed field excitations, yet experiments
have found no intrinsic constituent structure or finite form factor for
either. What matters for the present argument is that neither has a dense
tower of almost-degenerate internal states.\editorialnote{The lecture
associates the electron's anomalous magnetic moment with a measurable quantum
fuzz. In modern terminology the anomalous moment is a radiative dressing
effect, not evidence for a resolved finite electron radius. Pointlike
conventionally means that no intrinsic substructure or form factor has been
detected down to the tested scale.}}
\end{classroomqa}

\begin{classroomqa}
\audiencequestion{Does the size of a photon depend on its
frequency?\lecturetimestamp{00:36:49}}
\lecturerresponse{Its wavelength certainly does. That is distinct from an
intrinsic constituent size. Changing the frequency changes the wave packet
and wavelength, not a resolved internal radius of the photon.}
\end{classroomqa}

\section{Exact transverse mechanics at large momentum}

We now return to the relativistic difficulty. A rapidly rotating baton cannot
be treated nonrelativistically merely by putting its center of mass at rest.
Its ends still move relativistically relative to one another. The
nonrelativistic approximation fails because of internal relative motion.

There is, however, a different limit. Choose a \(z\)-axis and boost the entire
system to a very large positive longitudinal momentum \(P_z\). Keep its
transverse momentum

\begin{equation}
  p_\perp^2=p_x^2+p_y^2
\end{equation}

finite. The transverse projection then admits an exactly nonrelativistic
Hamiltonian form in the \(P_z\to\infty\) limit, even if the unboosted internal
motion was relativistic.

\begin{classroomqa}
\audiencequestion{Does infinite momentum mean that the massive system
actually reaches the speed of light?\lecturetimestamp{00:40:43}}
\lecturerresponse{No finite boost reaches it. One constructs quantities with
well-defined limits as \(P_z\) becomes arbitrarily large. The phrase names a
limiting procedure, not a finite physical state with \(v=c\).}
\end{classroomqa}

\begin{classroomqa}
\audiencequestion{Which axis must be chosen, and what if the object rotates
about or across that axis?\lecturetimestamp{00:41:14}}
\lecturerresponse{Any axis may be chosen and the whole system boosted along
it. What remains as explicit mechanics is the projection onto the
perpendicular plane. A rotating rod may look foreshortened or sweep through
that plane, but its transverse projection has the nonrelativistic form.}
\end{classroomqa}

Strings have an additional property that makes this representation unusually
powerful. In light-front gauge their longitudinal internal coordinates are
not independent oscillators. The constraints determine them from the
transverse motion. Knowing the transverse string is therefore enough to
reconstruct the physical degrees of freedom.

\begin{classroomqa}
\audiencequestion{Does the absence of an independent longitudinal coordinate
mean that the string is merely trapped in a thin
sandwich?\lecturetimestamp{00:42:26}}
\lecturerresponse{Lorentz contraction does make its longitudinal extent
small in the boosted frame, but the statement is stronger. Any residual
longitudinal motion is constrained by the transverse motion; it cannot be
specified independently.}
\end{classroomqa}

\begin{classroomqa}
\audiencequestion{Is this reduction to transverse data related to the
holographic principle?\lecturetimestamp{00:43:12}}
\lecturerresponse{It is closely connected to the line of thought that led to
holographic ideas, although pursuing that connection would take us away from
the present calculation. Here it serves only to justify working with the
transverse string.}
\end{classroomqa}

\subsection{Subtracting the boost and rescaling time}

Let \(M\) be the invariant mass of the complete system, including constituent
masses, kinetic energy, and binding energy. In units \(c=1\),

\begin{equation}
  E=\sqrt{P_z^2+p_\perp^2+M^2}.
  \label{eq:l02-relativistic-energy}
\end{equation}

For \(P_z\) much larger than \(p_\perp\) and \(M\),

\begin{equation}
  E
  =
  P_z+\frac{p_\perp^2+M^2}{2P_z}
  +O(P_z^{-3}).
  \label{eq:l02-large-pz}
\end{equation}

\begin{figure}[H]
  \centering
  \includegraphics[width=0.76\textwidth]{lecture_02_figure_08.png}
  \caption{The completed large-\(P_z\) expansion, with transverse momentum
  and invariant mass retained.}
  \lectureframe{00:47:30}
  \label{fig:l02-large-pz-board}
\end{figure}

The leading term is large, but it carries no information about the internal
state. In a process that conserves \(P_z\), such as a boosted atom changing
between internal levels without changing its total longitudinal momentum,
that term cancels from every energy difference. We therefore study

\begin{equation}
  E-P_z=\frac{p_\perp^2+M^2}{2P_z}+O(P_z^{-3}).
  \label{eq:l02-residual-energy}
\end{equation}

The denominator has a physical meaning. Since the Hamiltonian generates time
evolution, the small residual energy says that internal evolution appears
slow in laboratory time. This is time dilation. To retain the internal
motion, introduce a rescaled evolution parameter \(\tau\), or equivalently
multiply the residual Hamiltonian by \(P_z\):

\begin{equation}
  H_\perp
  \equiv
  \lim_{P_z\to\infty}P_z(E-P_z)
  =
  \frac12p_\perp^2+\frac12M^2.
  \label{eq:l02-transverse-hamiltonian}
\end{equation}

\begin{figure}[H]
  \centering
  \begin{minipage}[t]{0.48\textwidth}
    \centering
    \includegraphics[width=\textwidth]{lecture_02_figure_09.png}
    \par\smallskip\footnotesize (a) \(H=i\partial_t\) and the slowed internal
    evolution.
  \end{minipage}\hfill
  \begin{minipage}[t]{0.48\textwidth}
    \centering
    \includegraphics[width=\textwidth]{lecture_02_figure_10.png}
    \par\smallskip\footnotesize (b) The finite transverse Hamiltonian after
    time rescaling.
  \end{minipage}
  \caption{From the time-dilated residual energy to a finite
  two-dimensional Hamiltonian.}
  \lectureframe{00:49:20 and 00:52:40}
  \label{fig:l02-lightfront-board}
\end{figure}

This is precisely the form of a unit-mass particle moving in two dimensions:
kinetic energy \(p_\perp^2/2\) plus an internal energy \(M^2/2\). The unusual
feature is the square. In this description, the quantity playing the role of
rest or binding energy is proportional to invariant mass
squared.\editorialnote{With conventional light-front variables
\(P^\pm=(E\pm P_z)/\sqrt2\), the exact mass-shell relation is
\(2P^+P^-=p_\perp^2+M^2\). Equation
\eqref{eq:l02-transverse-hamiltonian} is the lecture's large-\(P_z\)
normalization of the same invariant relation.}

\begin{classroomqa}
\audiencequestion{What happened to the longitudinal momenta of the individual
parts of a composite system?\lecturetimestamp{00:51:30}}
\lecturerresponse{For a generic composite object one must retain the
fractions of total \(P_z\) carried by its constituents. The formula above
uses the total center-of-mass momentum. String theory is special because its
constraints determine the relative longitudinal motion from the transverse
degrees of freedom, so no independent longitudinal oscillator has to be
added.}
\end{classroomqa}

\section{A discrete open string}

The boosted string now moves explicitly only in the transverse plane. Label
its two coordinates \(X\) and \(Y\), and first approximate it by \(N\) mass
points joined by \(N-1\) springs. If \(\mu\) is the analog mass of each point
and \(k\) the nearest-neighbor spring constant, the energy is

\begin{equation}
  E_N
  =
  \sum_i
  \left[
    \frac{\mu}{2}
    \left(\dot X_i^2+\dot Y_i^2\right)
    +
    \frac{k}{2}
    \left((\Delta X_i)^2+(\Delta Y_i)^2\right)
  \right].
  \label{eq:l02-discrete-energy}
\end{equation}

Every \(X\)-term has an identical \(Y\)-term because the transverse plane has
two directions. The index \(i\) labels a point along the chain; it is not a
spatial coordinate.

\begin{figure}[H]
  \centering
  \begin{minipage}[t]{0.48\textwidth}
    \centering
    \includegraphics[width=\textwidth]{lecture_02_figure_11.png}
    \par\smallskip\footnotesize (a) The bead chain and its discrete energy.
  \end{minipage}\hfill
  \begin{minipage}[t]{0.48\textwidth}
    \centering
    \includegraphics[width=\textwidth]{lecture_02_figure_12.png}
    \par\smallskip\footnotesize (b) The large-\(N\) scaling of \(\mu\) and
    \(k\).
  \end{minipage}
  \caption{The regulated open string before taking the continuum limit.}
  \lectureframe{00:57:42 and 01:02:40}
  \label{fig:l02-chain-board}
\end{figure}

\begin{figure}[H]
  \centering
  \begin{tikzpicture}[>=Latex,thick]
    \foreach \i in {0,...,8} {
      \coordinate (p\i) at ({1.05*\i},{0.30*sin(55*\i)});
      \filldraw[white] (p\i) circle (0.10);
      \draw (p\i) circle (0.10);
    }
    \foreach \i in {0,...,7} {
      \pgfmathtruncatemacro{\j}{\i+1}
      \draw[decorate,decoration={zigzag,segment length=4pt,amplitude=2pt}]
        (p\i) -- (p\j);
    }
    \node[above=0.35cm of p0] {\(i=0\)};
    \node[above=0.35cm of p8] {\(i=N-1\)};
    \draw[<->] (0,-1.15) -- (8.4,-1.15)
      node[midway,below] {\(0\leq\sigma\leq\pi\)};
    \node at (4.2,1.25)
      {\(\mu=1/N,\quad k=N/\pi^2,\quad\Delta\sigma=\pi/N\)};
  \end{tikzpicture}
  \caption{A clean regulator picture. The points become lighter while the
  microscopic springs become stiffer as the chain is refined.}
  \label{fig:l02-chain-redraw}
\end{figure}

The parameter \(\sigma\) labels position along the chain. Choosing
\(0\leq\sigma\leq\pi\) is a convention, not a physical length measurement.
It reminds us that the string has two endpoints. Later a closed string will
be labeled periodically over \(0\leq\sigma<2\pi\).

The continuum limit must keep the collective mechanics finite. Fixed bead
mass would make the total analog mass diverge, so choose

\begin{equation}
  \mu=\frac{1}{N}.
  \label{eq:l02-mu-scaling}
\end{equation}

Then \(N\mu=1\). This is not the relativistic rest mass \(M\); it is the
inertial coefficient in the transverse analog model.

Springs in series become softer as their number increases. To preserve a
finite stiffness for the whole chain, each microscopic spring must become
stiffer:

\begin{equation}
  k=\frac{N}{\pi^2}
  =\frac{1}{\pi\,\Delta\sigma}.
  \label{eq:l02-k-scaling}
\end{equation}

Only \(k\propto N\) is essential. The factor of \(\pi^2\) is a normalization
chosen to remove stray factors of \(\pi\) from the continuum formulas.

\begin{classroomqa}
\audiencequestion{Is \(\mu\) the physical rest mass of each string
point?\lecturetimestamp{01:10:34}}
\lecturerresponse{No. It is the nonrelativistic analog mass appearing in the
transverse chain. Calling it \(M\) would confuse it with the invariant rest
mass of the entire relativistic string state. Its only present role is to
produce the correct continuum inertia.}
\end{classroomqa}

\section{The continuum wave system}

Apply the rules prepared at the beginning. Since
\(1/N=\Delta\sigma/\pi\),

\begin{align}
  \sum_i\frac{\mu}{2}\dot X_i^2
  &=
  \frac{1}{2\pi}
  \sum_i\dot X_i^2\,\Delta\sigma
  \longrightarrow
  \frac{1}{2\pi}
  \int_0^\pi
  \left(\partial_\tau X\right)^2d\sigma ,
  \label{eq:l02-kinetic-limit}\\
  \sum_i\frac{k}{2}(\Delta X_i)^2
  &\longrightarrow
  \frac{1}{2\pi}
  \int_0^\pi
  \left(\partial_\sigma X\right)^2d\sigma .
  \label{eq:l02-potential-limit}
\end{align}

For the second line, use
\(\Delta X_i=\partial_\sigma X\,\Delta\sigma\) together with
\(k=1/(\pi\Delta\sigma)\). The same calculation applies to \(Y\).

\begin{figure}[H]
  \centering
  \begin{minipage}[t]{0.48\textwidth}
    \centering
    \includegraphics[width=\textwidth]{lecture_02_figure_13.png}
    \par\smallskip\footnotesize (a) Sum-to-integral and
    difference-to-derivative replacements.
  \end{minipage}\hfill
  \begin{minipage}[t]{0.48\textwidth}
    \centering
    \includegraphics[width=\textwidth]{lecture_02_figure_14.png}
    \par\smallskip\footnotesize (b) The completed continuum energy.
  \end{minipage}
  \caption{The discrete chain becoming a field on the \(\sigma\)-interval.}
  \lectureframe{01:05:40 and 01:07:00}
  \label{fig:l02-continuum-board}
\end{figure}

The resulting energy and Lagrangian are

\begin{align}
  E
  &=
  \frac{1}{2\pi}
  \int_0^\pi d\sigma\,
  \left[
    (\partial_\tau X)^2+(\partial_\tau Y)^2
    +(\partial_\sigma X)^2+(\partial_\sigma Y)^2
  \right],
  \label{eq:l02-continuum-energy}\\
  L
  &=
  \frac{1}{2\pi}
  \int_0^\pi d\sigma\,
  \left[
    (\partial_\tau X)^2+(\partial_\tau Y)^2
    -(\partial_\sigma X)^2-(\partial_\sigma Y)^2
  \right].
  \label{eq:l02-continuum-lagrangian}
\end{align}

\begin{figure}[H]
  \centering
  \includegraphics[width=0.78\textwidth]{lecture_02_figure_15.png}
  \caption{The continuum Lagrangian with kinetic and stretching terms
  carrying opposite signs.}
  \lectureframe{01:08:48}
  \label{fig:l02-lagrangian-board}
\end{figure}

If we rename \(\sigma\) as an ordinary spatial coordinate and \(X\) as a
field \(\phi\), these are the energy and Lagrangian of a one-dimensional wave
field. Varying \eqref{eq:l02-continuum-lagrangian} gives

\begin{equation}
  \left(\partial_\tau^2-\partial_\sigma^2\right)X=0,
  \qquad
  \left(\partial_\tau^2-\partial_\sigma^2\right)Y=0.
  \label{eq:l02-wave-equation}
\end{equation}

Nothing in the immediate argument requires us to solve these equations
directly. The useful picture is simply that waves travel left and right along
the string and reflect at the endpoints. A Dirichlet reflection reverses the
sign of the displacement; a Neumann reflection preserves it.

\begin{classroomqa}
\audiencequestion{What exactly is the time variable
\(\tau\)?\lecturetimestamp{01:08:54}}
\lecturerresponse{It is the rescaled time used to follow the internal motion
in the large-\(P_z\) description. A clock boosted with the system runs slowly
relative to laboratory time; \(\tau\) removes that dilation so the internal
oscillations retain a finite rate. For this transverse mechanics it can be
treated simply as time.\editorialnote{The lecture informally calls \(\tau\)
proper time. For an extended string, the precise light-front evolution
parameter is a gauge-fixed rescaled time rather than a single invariant proper
time shared by every point. The intended time-dilation argument is
unchanged.}}
\end{classroomqa}

When the center of mass is at rest in the transverse plane, the internal
part of \eqref{eq:l02-continuum-energy} is, in the normalization of
\eqref{eq:l02-transverse-hamiltonian}, proportional to \(M^2\). That is why
quantizing this apparently ordinary wave system will reveal the relativistic
mass spectrum.

\begin{classroomqa}
\audiencequestion{Does \(P_z(E-P_z)\) have a
name?\lecturetimestamp{01:11:00}}
\lecturerresponse{It is commonly called a light-cone or light-front energy.
Light front better describes the null plane used as the evolution surface;
the traditional terminology is not important for the calculation.}
\end{classroomqa}

\begin{classroomqa}
\audiencequestion{Is the word cone referring to the usual spacetime light
cone, and did Dirac introduce this terminology?\lecturetimestamp{01:11:17}}
\lecturerresponse{The null-coordinate construction can be drawn with
45-degree axes, but the evolution surface here is a light front rather than
the full expanding cone. Dirac developed related forms of relativistic
dynamics; the historical naming question was left open and does not affect
the equations.}
\end{classroomqa}

\section{Why a free endpoint is Neumann}

We can now derive rather than postulate the endpoint condition. Newton's law
is legitimate because the boosted transverse mechanics has exactly the
nonrelativistic form. An interior mass point is pulled by springs on both
sides. The final point is pulled only by the spring joining it to its
neighbor.

\begin{classroomqa}
\audiencequestion{Why are we allowed to use Newton's law while discussing a
relativistic string?\lecturetimestamp{01:14:52}}
\lecturerresponse{We are not applying Newtonian mechanics to the full
unboosted spacetime motion. We are applying it to the transverse variables
in the infinite-momentum description, whose Hamiltonian has exactly the
nonrelativistic form.}
\end{classroomqa}

At the right endpoint, Hooke's law gives, up to orientation,

\begin{equation}
  F_{\mathrm{end}}
  =
  -k\,\Delta X
  =
  -\frac{1}{\pi\Delta\sigma}
   \left(\partial_\sigma X\,\Delta\sigma\right)
  =
  -\frac{1}{\pi}\partial_\sigma X .
  \label{eq:l02-endpoint-force}
\end{equation}

\begin{figure}[H]
  \centering
  \begin{minipage}[t]{0.48\textwidth}
    \centering
    \includegraphics[width=\textwidth]{lecture_02_figure_16.png}
    \par\smallskip\footnotesize (a) Endpoint force proportional to the local
    slope.
  \end{minipage}\hfill
  \begin{minipage}[t]{0.48\textwidth}
    \centering
    \includegraphics[width=\textwidth]{lecture_02_figure_17.png}
    \par\smallskip\footnotesize (b) The finite-acceleration condition.
  \end{minipage}
  \caption{The Newton-law derivation of the free-end boundary condition.}
  \lectureframe{01:18:20 and 01:20:10}
  \label{fig:l02-neumann-board}
\end{figure}

Newton's equation for the final point is

\begin{equation}
  \mu\,\partial_\tau^2X=F_{\mathrm{end}},
  \qquad
  \mu=\frac1N=\frac{\Delta\sigma}{\pi}.
  \label{eq:l02-endpoint-newton}
\end{equation}

Combining \eqref{eq:l02-endpoint-force} and
\eqref{eq:l02-endpoint-newton},

\begin{equation}
  \partial_\sigma X
  =
  -\Delta\sigma\,\partial_\tau^2X
  \quad\text{at the right endpoint},
  \label{eq:l02-endpoint-scaling}
\end{equation}

with the force sign reversed at the left endpoint. If
\(\partial_\sigma X\) approached a nonzero constant while
\(\Delta\sigma\to0\), the endpoint acceleration would diverge like \(N\).
A sensible continuum limit with finite acceleration therefore requires

\begin{equation}
  \left.\partial_\sigma X\right|_{0,\pi}=0,
  \qquad
  \left.\partial_\sigma Y\right|_{0,\pi}=0.
  \label{eq:l02-open-neumann}
\end{equation}

\begin{figure}[H]
  \centering
  \begin{tikzpicture}[>=Latex,thick,scale=1.0]
    \coordinate (a) at (0,0);
    \coordinate (b) at (4.5,0.75);
    \filldraw[white] (a) circle (0.12);
    \draw (a) circle (0.12);
    \filldraw[white] (b) circle (0.12);
    \draw (b) circle (0.12);
    \draw[decorate,decoration={zigzag,segment length=6pt,amplitude=3pt}]
      (a) -- (b);
    \draw[dashed] (a) -- (0,-1.2);
    \draw[dashed] (b) -- (4.5,-1.2);
    \draw[<->] (4.85,0) -- (4.85,0.75)
      node[midway,right] {\(\Delta X\)};
    \draw[->,red!70!black] (b) -- (5.75,0.95)
      node[right] {\(F_{\rm end}=-k\Delta X\)};
    \node[below] at (0,-1.2) {\(N-1\)};
    \node[below] at (4.5,-1.2) {\(N\)};
    \node at (2.5,-2.0)
      {\(\mu\ddot X=-k\Delta X
      \ \Longrightarrow\
      \partial_\sigma X=-\Delta\sigma\,\ddot X\to0\)};
  \end{tikzpicture}
  \caption{The last point has vanishing mass in the regulator limit. A finite
  acceleration is possible only when the endpoint slope vanishes.}
  \label{fig:l02-neumann-redraw}
\end{figure}

The same argument at both ends and in both transverse directions produces
Neumann boundary conditions for a free open string.

\section{Normal modes of the open string}

We now have a complete classical system: a Lagrangian and its boundary
conditions. Quantization becomes simple after Fourier analysis. Because the
endpoints are Neumann,

\begin{align}
  X(\sigma,\tau)
  &=
  \sum_{n=0}^{\infty}X_n(\tau)\cos(n\sigma),
  \label{eq:l02-X-expansion}\\
  Y(\sigma,\tau)
  &=
  \sum_{n=0}^{\infty}Y_n(\tau)\cos(n\sigma).
  \label{eq:l02-Y-expansion}
\end{align}

The cosines contain all \(\sigma\)-dependence. The coefficients contain all
time-dependence; the string wiggles because \(X_n(\tau)\) and \(Y_n(\tau)\)
change.

\begin{figure}[H]
  \centering
  \includegraphics[width=0.78\textwidth]{lecture_02_figure_18.png}
  \caption{The Neumann mode expansion placed beside the orthogonality rule
  that will collapse the double sums.}
  \lectureframe{01:21:45}
  \label{fig:l02-mode-expansion-board}
\end{figure}

\subsection{Kinetic term and the zero mode}

Time differentiation gives

\begin{equation}
  \partial_\tau X
  =
  \sum_{n=0}^{\infty}\dot X_n(\tau)\cos(n\sigma).
  \label{eq:l02-time-mode}
\end{equation}

Squaring produces a double sum over \(n\) and \(m\), but
\eqref{eq:l02-cos-orthogonality} removes every term with \(n\neq m\).
Treating the constant mode separately,

\begin{equation}
  \frac{1}{2\pi}\int_0^\pi
  (\partial_\tau X)^2\,d\sigma
  =
  \frac12\dot X_0^{\,2}
  +
  \frac14\sum_{n=1}^{\infty}\dot X_n^{\,2}.
  \label{eq:l02-mode-kinetic}
\end{equation}

\begin{figure}[H]
  \centering
  \includegraphics[width=0.76\textwidth]{lecture_02_figure_19.png}
  \caption{The orthogonality integral reducing the kinetic double sum to the
  zero mode plus a diagonal sum.}
  \lectureframe{01:29:10}
  \label{fig:l02-mode-kinetic-board}
\end{figure}

The constant coefficient is the average position,

\begin{equation}
  X_0(\tau)=\frac1\pi\int_0^\pi X(\sigma,\tau)\,d\sigma ,
  \label{eq:l02-center-of-mass}
\end{equation}

and similarly for \(Y_0\). With uniform analog mass density, these are the
center-of-mass coordinates. They have kinetic energy but no restoring force:
the entire string is free to move across the transverse plane.

\subsection{Stretching term and oscillator frequencies}

Spatial differentiation gives

\begin{equation}
  \partial_\sigma X
  =
  -\sum_{n=1}^{\infty}nX_n(\tau)\sin(n\sigma).
  \label{eq:l02-spatial-mode}
\end{equation}

The zero mode disappears because a uniform translation does not stretch the
string. Using sine orthogonality,

\begin{equation}
  \frac{1}{2\pi}\int_0^\pi
  (\partial_\sigma X)^2\,d\sigma
  =
  \frac14\sum_{n=1}^{\infty}n^2X_n^{\,2}.
  \label{eq:l02-mode-potential}
\end{equation}

\begin{figure}[H]
  \centering
  \includegraphics[width=0.76\textwidth]{lecture_02_figure_20.png}
  \caption{Differentiating the cosine series produces the sine double sum;
  orthogonality again enforces \(n=m\).}
  \lectureframe{01:31:40}
  \label{fig:l02-mode-potential-board}
\end{figure}

The \(X\)-sector Lagrangian is therefore

\begin{equation}
  L_X
  =
  \frac12\dot X_0^{\,2}
  +
  \frac14\sum_{n=1}^{\infty}
  \left(
    \dot X_n^{\,2}-n^2X_n^{\,2}
  \right),
  \label{eq:l02-X-lagrangian}
\end{equation}

with an identical expression for \(Y\). The corresponding energy has a plus
sign in front of \(n^2X_n^2\). Every nonzero mode is an independent harmonic
oscillator, and comparison with \eqref{eq:l02-oscillator-lagrangian} gives

\begin{equation}
  \omega_n=n,
  \qquad n=1,2,3,\ldots .
  \label{eq:l02-mode-frequency}
\end{equation}

\begin{figure}[H]
  \centering
  \includegraphics[width=0.80\textwidth]{lecture_02_figure_21.png}
  \caption{The decoupled oscillator Lagrangian and the conclusion
  \(\omega_n=n\), with the identical \(Y\)-tower indicated below.}
  \lectureframe{01:35:50}
  \label{fig:l02-oscillator-modes-board}
\end{figure}

The normal modes are simply the harmonics of the string. The first bends
slowly, the second has twice the frequency, the third three times, and so on.
Orthogonality is what prevents one harmonic from coupling to another in the
free theory.

\begin{classroomqa}
\audiencequestion{What happened to the \(n=0\)
oscillator?\lecturetimestamp{01:34:49}}
\lecturerresponse{It has zero frequency because it is not an internal
oscillator at all. It translates the center of mass. No spring restores the
whole string to a preferred transverse position, so the zero mode is a free
particle coordinate.}
\end{classroomqa}

Dropping the center-of-mass kinetic term leaves the internal energy

\begin{equation}
  E_{\mathrm{int}}
  =
  \frac14
  \sum_{n=1}^{\infty}
  \left[
    \dot X_n^{\,2}+n^2X_n^{\,2}
    +\dot Y_n^{\,2}+n^2Y_n^{\,2}
  \right].
  \label{eq:l02-internal-energy}
\end{equation}

In the light-front normalization this is proportional to the invariant
mass squared of the string state. We have thus converted the relativistic
mass problem into the quantum mechanics of two countable oscillator towers.

\section{Quantized levels and particle behavior}

Let \(a_n^\dagger\) create one quantum in the \(X\)-oscillator of frequency
\(n\), and \(b_n^\dagger\) do the same for \(Y\). The lecture postpones a
systematic operator treatment, but this notation makes the level counting
transparent. Relative to the ground-state offset,

\begin{equation}
  \Delta E_{\mathrm{int}}
  =
  \sum_{n=1}^{\infty}
  n\left(Q_n^X+Q_n^Y\right),
  \qquad
  Q_n^X,Q_n^Y\in\{0,1,2,\ldots\},
  \label{eq:l02-quantized-energy}
\end{equation}

where \(\hbar=1\). The integer \(n\) labels which oscillator is excited;
\(Q_n\) counts how many quanta occupy it.\editorialnote{Each oscillator also
contributes a zero-point energy. The
infinite sum requires regularization and supplies the intercept of the string
mass formula. The lecture deliberately leaves the ground-state offset
unspecified here and counts only excitation levels.}

The oscillator vacuum \(\ket{0}\) is not empty spacetime. It is one open
string with no internal oscillator excited. It has energy, momentum, and a
definite place in the string spectrum.

At level one there are two states,

\begin{equation}
  a_1^\dagger\ket{0},
  \qquad
  b_1^\dagger\ket{0},
  \label{eq:l02-level-one}
\end{equation}

corresponding to the lowest harmonic polarized along \(X\) or \(Y\).
Linear combinations describe any transverse polarization.

At level two there are five states:

\begin{equation}
  (a_1^\dagger)^2\ket{0},
  \quad
  a_1^\dagger b_1^\dagger\ket{0},
  \quad
  (b_1^\dagger)^2\ket{0},
  \quad
  a_2^\dagger\ket{0},
  \quad
  b_2^\dagger\ket{0}.
  \label{eq:l02-level-two}
\end{equation}

Three use two quanta of the first harmonic; two use one quantum of the second
harmonic. Both constructions cost two units.

\begin{figure}[H]
  \centering
  \includegraphics[width=0.76\textwidth]{lecture_02_figure_22.png}
  \caption{The five level-two states being recounted after an audience
  challenge: two double excitations, one mixed excitation, and two
  second-harmonic excitations.}
  \lectureframe{01:44:00}
  \label{fig:l02-levels-board}
\end{figure}

\begin{figure}[H]
  \centering
  \begin{tikzpicture}[>=Latex,thick]
    \draw[->] (0,0) -- (0,5.1) node[above] {\(\Delta E_{\rm int}\)};
    \draw[blue!70!black] (-0.65,0.55) -- (0.65,0.55);
    \node[right] at (0.85,0.55) {\(\ket{0}\) \quad multiplicity \(1\)};
    \draw[blue!70!black] (-0.65,2.15) -- (0.65,2.15);
    \node[right] at (0.85,2.15)
      {\(a_1^\dagger\ket{0},\ b_1^\dagger\ket{0}\)
      \quad multiplicity \(2\)};
    \draw[blue!70!black] (-0.65,3.75) -- (0.65,3.75);
    \node[right,text width=8.6cm] at (0.85,3.75)
      {\((a_1^\dagger)^2\ket{0},
      a_1^\dagger b_1^\dagger\ket{0},
      (b_1^\dagger)^2\ket{0},
      a_2^\dagger\ket{0},\ b_2^\dagger\ket{0}\)
      \quad multiplicity \(5\)};
    \node[left] at (-0.75,0.55) {\(0\)};
    \node[left] at (-0.75,2.15) {\(1\)};
    \node[left] at (-0.75,3.75) {\(2\)};
  \end{tikzpicture}
  \caption{The first three excitation levels of the two-transverse-direction
  open string.}
  \label{fig:l02-levels-redraw}
\end{figure}

\begin{classroomqa}
\audiencequestion{Are there four or five states at the second excited
level?\lecturetimestamp{01:43:04}}
\lecturerresponse{Five. Exciting \(X_1\) twice, exciting \(Y_1\) twice, and
exciting each once give three states. Exciting \(X_2\) once or \(Y_2\) once
gives two more. The proposed exercise is to count the next two levels.}
\end{classroomqa}

The essential conclusion is already visible. Successive levels are separated
by finite gaps, and each has a finite multiplicity. The quantized string is
therefore particle-like by the spectral criterion that distinguished it from
coffee, gas, or an ordinary violin string. Later we shall determine which
states are physically welcome and which pathologies must be removed.

\section{Why open-string interactions produce closed strings}

An open string has two endpoints. Its elementary interactions join and split
strings. In the hadronic picture one may imagine a quark at one end and an
antiquark at the other; their annihilation joins two string segments. The
fundamental strings of interest later need not carry that literal hadronic
endpoint interpretation, but the local joining rule remains.

Now let one open string fluctuate until its own endpoints approach one
another. Locally, an endpoint cannot know that the other endpoint belongs to
the same string rather than a different one. If endpoints may join at all,
they may join here, producing a closed loop.

\begin{figure}[H]
  \centering
  \includegraphics[width=0.70\textwidth]{lecture_02_figure_23.png}
  \caption{The board sketch of an open string fluctuating until its two
  endpoints meet and form a closed loop.}
  \lectureframe{01:47:00}
  \label{fig:l02-closed-board}
\end{figure}

\begin{figure}[H]
  \centering
  \begin{tikzpicture}[>=Latex,thick]
    \begin{scope}
      \draw[blue!70!black]
        (0,0) .. controls (0.8,1.2) and (1.7,-1.0) .. (2.6,0.45)
        .. controls (3.4,1.55) and (4.0,0.3) .. (4.7,0.85);
      \fill (0,0) circle (0.09);
      \fill (4.7,0.85) circle (0.09);
      \node at (2.35,-1.25) {open string fluctuation};
    \end{scope}
    \draw[->,line width=1pt] (5.25,0.3) -- (6.65,0.3)
      node[midway,above] {join};
    \begin{scope}[xshift=8.1cm]
      \draw[orange!85!black]
        (0,0) .. controls (0.3,1.5) and (2.5,1.8) .. (3.1,0.3)
        .. controls (3.5,-1.1) and (0.7,-1.5) .. (0,0);
      \node at (1.55,-1.85) {closed string};
    \end{scope}
  \end{tikzpicture}
  \caption{Once endpoint joining is allowed, an open string can close on
  itself.}
  \label{fig:l02-closed-redraw}
\end{figure}

Thus a theory may contain closed strings without open strings, but an
interacting theory with open strings cannot consistently avoid closed
strings. This asymmetry is the important prediction. The spectral
identifications will be developed later, but the destination is already
visible: massless open-string states often supply gauge particles such as
photons, while the closed-string spectrum contains a massless spin-two state,
the graviton. Closed strings, and therefore gravity, are not an optional
decoration once string theory is interacting.
